\section{Methods and theory}
\label{sec:methods}
% Introduction to methods chosen here.
\subsection{Rigorous Coupled Wave Analysis}
\label{classical_solvers}
This subsection will introduce the concept of Rigorous Coupled Wave Analysis (RCWA), a method in numerical photonics suited well
for stacked periodic structures in the x and y axis with a period $\Lambda_x,\Lambda_y$ respectively, and where each layer is homogeneous in the z-direction\cite{torcwa}.
The structure's optical properties are completely described by the relative permitivity $\epsilon$ and relative permeability $\mu$ of each layer, meaning that each structure
will have a geometry described by
\begin{equation}
\label{eq:permittivity_periodicity}
\epsilon(x,y,z)_N = \epsilon(x + m \cdot \Lambda_x, y + n \cdot \Lambda_y)_N, m,n \in \mathbb{Z}
\end{equation}
and 
\begin{equation}
\label{eq:permeability_periodicity}
\mu(x,y,z)_N = \mu(x + m \cdot \Lambda_x, y + n \cdot \Lambda_y)_N, m,n \in \mathbb{Z},
\end{equation}
where $m$ and $n$ are integers and the subscript $N$ denotes the $N$'th layer of the structure. For the entirety of this project, the designed
fourier patterned light trapping layers will be designed using crystalline silicon ($Si$), Titanium Dioxide ($TiO_2$) and Silicon Nitride ($Si_3N_4$), which
are all nonmagnetic, so the relative permeability $\mu = 1$. In general, $\epsilon$ and $\mu$ will be $3 \times 3$ tensors, but for isotropic and linear materials this is 
a bit simpler. The equations solved for are of course the Maxwell equations, namely Faraday's law
\begin{equation}
\label{eq:faraday}
\nabla \times \mathbf{E} = -\mu_0\mu(x,y)\frac{\partial\mathbf{H}}{\partial t}
\end{equation}
and Ampere's law
\begin{equation}
\label{eq:ampere}
\nabla \times \mathbf{H} = \epsilon_0\epsilon(x,y)\frac{\partial\mathbf{E}}{\partial t},
\end{equation}
Where $\epsilon_0$ and $\mu_0$ are the vacuum permitivity and vacuum permeability respectively. The structures are assumed free of charge and current,
$$
\nabla \cdot (\epsilon \mathbf{E}) = 0
$$
and
$$
\nabla \cdot (\mu \mathbf{H}) = 0.
$$
At a single frequency, the time dependency of Maxwell's equations may be solved by plane waves, see \cite{griffiths} section 9.2.1 for example,
meaning
\begin{equation}
\label{eq:planewaveE}
\mathbf{E} = \mathbf{E}_0 \exp{i(\mathbf{k}\cdot\mathbf{r} - \omega t)}
\end{equation}
and
\begin{equation}
\label{eq:planewaveH}
\mathbf{H} = \mathbf{H}_0 \exp{i(\mathbf{k}\cdot\mathbf{r} - \omega t)},
\end{equation}
Where $\mathbf{E}_0$ and $\mathbf{H}_0$ are the (complex) amplitudes of the fields, with the physical fields being the real parts of
$\mathbf{E}$ and $\mathbf{H}$, $\mathbf{k}$ being the wavevector, i.e. the spatial frequency, and $\omega$ being the frequency of the fields, also given by the dispersion relation
$\omega = \frac{c|\mathbf{k}|}{n(\lambda)}$, where $\lambda$ is the wavelength of the light.

Inserting equations \ref{eq:planewaveE} and \ref{eq:planewaveH} in equations \ref{eq:faraday} and \ref{eq:ampere} converts the time derivative to
an algebraic factor of $-i\omega t$, meaning that the equations appear in their time harmonic form, common to many frequency domain numerical methods in photonics,
\begin{equation}
\label{eq:THfaraday}
\nabla \times \mathbf{E} = i\omega\mu_0\mu(x,y)\mathbf{H}
\end{equation}
\begin{equation}
\label{eq:THampere}
\nabla \times \mathbf{H} = -i\omega\mu_0\mu(x,y)\mathbf{E}.
\end{equation}
Owing to the periodicity of the system, see equations \ref{eq:permittivity_periodicity} and  \ref{eq:permeability_periodicity}, both the electromagnetic fields and the permitivity and permeability may
be written as fourier series in the x- and y-directions as below,
\begin{equation}
\label{eq:fourierPermitivity}
\epsilon(x,y) = \sum_{m=-M}^{M}\sum_{n=-N}^{N}\epsilon_{m,n}\exp{i(mL_xx+nL_yy)},
\end{equation}
\begin{equation}
\label{eq:fourierPermeability}
\mu(x,y) = \sum_{m=-M}^{M}\sum_{n=-N}^{N}\mu_{m,n}\exp{i(mL_xx+nL_yy)},
\end{equation}
\begin{equation}
\label{eq:fourierFaraday}
\mathbf{E} = \exp{i\mathbf{k}\cdot\mathbf{r}}\sum_{m=-M}^{M}\sum_{n=-N}^{N}\mathbf{E}_{m,n}\exp{i(mL_xx+nL_yy)},
\end{equation}
\begin{equation}
\label{eq:fourierAmpere}
\mathbf{H} = \exp{i\mathbf{k}\cdot\mathbf{r}}\sum_{m=-M}^{M}\sum_{n=-N}^{N}\mathbf{H}_{m,n}\exp{i(mL_xx+nL_yy)},
\end{equation}
where $L_x = \frac{2\pi}{\Lambda_x}$ and $L_y = \frac{2\pi}{\Lambda_y}$, 
and the parameters $\epsilon_{m,n}, \mu_{m,n}, \mathbf{E}_{m,n}$ and $\mathbf{H}_{m,n}$ are the fourier amplitudes of the $m,n$'th diffraction order.



% Detail the rigorous coupled-wave analysis (Torcwa) and any classical optimization approaches used for data generation or baseline comparisons.

\subsection{Machine Learning Surrogates}
\label{subsec:ml_surrogates}
% Describe the forward models (e.g., MLP, SpatialCNN, SIREN) and the inverse/generative models (e.g., Tandem networks, Contrastive VAE).

\subsection{Optimization Strategies}
\label{subsec:optimization_strategies}
% Outline the active learning loop, surrogate optimization techniques, and how candidates are proposed and verified.
